\documentclass[11pt]{article}
\usepackage[letterpaper,margin=1.1in]{geometry}
\usepackage{mathptmx}
\usepackage[T1]{fontenc}
\usepackage{microtype}
\usepackage{booktabs}
\usepackage{enumitem}
\usepackage[hidelinks]{hyperref}
\setlength{\parindent}{0pt}
\setlength{\parskip}{0.6em}
\setlist[itemize]{itemsep=0.1em,topsep=0.2em}
\setlist[enumerate]{itemsep=0.1em,topsep=0.2em}
\newcommand{\sect}[1]{\vspace{0.9em}{\large\textbf{#1}}\par\vspace{-0.2em}}

\begin{document}

\noindent Linguistics 150C \hfill \textbf{Final paper guide}\\[-0.6em]
\noindent\rule{\linewidth}{0.4pt}

For the final paper, you will carry out independent research on one language, country, or island in the Pacific region. Your topic can be anything related to one or more of the course readings, as long as I approve it. You will develop your topic into a paper of at least 2,500 words that cites at least 7 academic sources. The final paper replaces a final exam. Altogether, the final paper counts for 25\% of your course grade, and the final presentation counts for another 15\%. The paper is submitted in four stages, plus a reflection, so that I can give you regular feedback. Below are the due dates and each component's share of the paper grade and of your overall course grade:

\begin{center}
\begin{tabular}{@{}llrr@{}}
\toprule
 & \textbf{Due date} & \textbf{Of paper} & \textbf{Of course}\\
\midrule
Topic idea & Friday, October 30 & 5\% & 1.25\%\\
Preliminary sources and research question & Friday, November 6 & 15\% & 3.75\%\\
Draft & Friday, November 20 & 25\% & 6.25\%\\
Final paper & Tuesday, December 15 & 45\% & 11.25\%\\
Final reflection & Tuesday, December 15 & 10\% & 2.5\%\\
\bottomrule
\end{tabular}
\end{center}

The presentation takes place between the draft and the final paper (December 1--10) and counts for 15\% of your course grade. Submit all written components through Lamākū. Late components lose 5\% for each day after the deadline (syllabus, Section 3.1).

\sect{Topic idea}
In a brief statement, identify a topic or a few topics that you would like to pursue in your research, and say why. Your topic must be about a language, country, or island in the Pacific region, and it should connect to something we have read or discussed in the course (for example, multilingualism, pidgins and creoles, borrowing, language contact, language endangerment, revitalization, or documentation). You can be as general or as specific as you like here; the idea is to start thinking seriously about where you want to go. If you have more than one idea, list them all, and tell me which one you are leaning toward. For this statement I am envisioning a length of 100--150 words, but a longer statement is fine.

After reading your topics, I will let you know whether they work and suggest directions to take, including readings to look for.

\sect{Preliminary sources and initial research question}
Find at least three academic sources that you think might be appropriate to your topic. Academic sources are scholarly books, book chapters, and peer-reviewed articles, including reference grammars and dictionaries for the language you choose. For each source, give a full reference and a few sentences about why it fits your topic: what it covers, and what you expect to take from it. Then state your initial research question (or questions), and write a short explanation of what you anticipate you will learn by answering it. Your question does not have to be final, and it will probably change as you read, but it should be specific enough to guide your reading. For example, ``How did the Spanish colonial period shape the vocabulary of Chamorro?'' is a question; ``Chamorro'' is a topic.

Use any citation style you like (APA, MLA, Unified Style Sheet for Linguistics, etc.), but be consistent here and in the final paper. I will respond with feedback, additional sources, and ideas for narrowing or reshaping your question.

\sect{Draft}
Write a complete first draft of your paper, at least 1,750 words (excluding references). The draft does not have to be polished, but it does need to be complete: it should have an introduction, a body, and a conclusion, and it should cite sources as you go. A strong draft will:

\begin{itemize}
\item Introduce your language, country, or island and the phenomenon you plan to discuss, and state your research question clearly
\item Present an argument or answer to your question that is supported by evidence from academic sources, with in-text citations
\item Connect your topic explicitly to at least one idea or reading from the course
\item Have a clear structure, with body paragraphs that each support one point and point back to your research question
\item Include a reference list in a consistent citation style
\end{itemize}

I will read your draft and return feedback before the final paper is due. If you have written fewer than 1,750 words, submit what you have; part of the draft is better than none.

\sect{Presentation}
You will present your project during Weeks 15 and 16 (December 1, 3, 8, and 10). Each student has a 10-minute slot: 7--8 minutes of presentation and 2--3 minutes of questions. You will sign up for a slot on a Google Sheet in Week 14. The presentation should be given with slides, and it is graded on:

\begin{itemize}
\item A clear research topic and research question
\item Accurate and relevant content
\item Use of at least three academic sources
\item Quality and clarity of the slides
\item Presentation flow and organization
\end{itemize}

A good presentation introduces your language and topic, states your research question, summarizes what the sources say, and shares your answer or current thinking. You can also use the presentation to ask your classmates for help with open problems. The presentation comes after the draft and before the final paper, so use the questions you get to improve your paper.

\sect{Final paper}
Your final paper should be a complete, revised version of your draft that takes into account my feedback and the feedback from your presentation. It must be at least 2,500 words, excluding the reference list, and must cite at least 7 academic sources, both in the text and in a reference list, in one consistent citation style. More important than the word count: the paper should be an academic essay with a clear topic, a clear argument, supporting evidence, and a conclusion, not a collection of facts. Please structure your paper as follows:

\begin{enumerate}[label=\alph*.]
\item \textbf{Introduction.} Give basic information about your language, country, or island and the phenomenon you plan to discuss. State your research question and your answer to it (your thesis).
\item \textbf{Background and analysis.} Take the reader through the evidence in organized sections or paragraphs. Each paragraph should make one point, support it with evidence from your sources, and show how it relates to your question. Connect your topic to the course readings and concepts.
\item \textbf{Conclusion.} Summarize your main points without copying them, say what further research could or should be done, and end on a clear statement about your language or about language in general that your paper supports.
\item \textbf{References.} List all of the sources you cite, alphabetically, in a consistent citation style.
\end{enumerate}

The final paper is due on Tuesday, December 15, by 11:45 am (our final exam period). See Section 3.3 of the syllabus for the course's generative AI policy and Section 3.2 for academic honesty.

\sect{Final reflection}
Submitted at the same time as your final paper, the reflection is a short piece of about 300 words about the process of researching and writing it. Write in the first person. You might consider questions such as: How did your topic and question change from the topic idea to the final paper? What was the most difficult or surprising part of the research? How did feedback on your sources, draft, and presentation change your paper? What would you do differently next time, and what did you learn about the language or about doing linguistic research?

\sect{Optional one-on-one meeting (extra credit)}
In Week 14 you can earn up to 4\% extra credit by meeting with me for 7 minutes, by Zoom or in person at my office hours, to discuss your plans for the paper and the presentation. You must have submitted your initial research question on time and come prepared to discuss your topic and presentation ideas. I am also happy to meet at any other point in the semester; come to office hours (Wednesdays and Fridays, 1:30--2:30, Moore 574) or email me at klineal@hawaii.edu.

\end{document}
